AutoResearch addresses a practical bottleneck in agent-assisted research: a fluent manuscript can hide whether claims, evidence, protocols, reviews, and responses are actually connected. Reporting and peer-review policies already expect inspectable claims, reproducible evidence, and transparent limitations \cite{nature-reporting,ieee-policies}. The system therefore treats a research node as the unit of work. Each node has a path, prompt assets, local artifacts, review status, and explicit closure conditions. This design makes a paper draft inspectable as a chain of bounded research decisions rather than a single prose surface.

The draft makes a narrower process contribution: it specifies that substantive manuscript claims should point to evidence, evidence items should retain explicit boundaries, and phase transitions should be backed by review evidence or recorded gate logic. In the current repository state, that contribution is supported as a process and documentation claim. It is not yet supported as a final real-data performance claim.

The current TeX synchronization carries forward only the evidence already accepted by upstream nodes. Earlier P1 nodes established data provenance, pseudocode, repository blueprinting, bounded synthetic/offline validation, result-state separation, and a draft result figure. The only quantitative signal allowed into this draft is the P1\_09 figure package derived from P1\_04/P1\_05 synthetic/offline evidence. The introduction therefore frames the paper around auditability and research-operation discipline, while reserving stronger empirical claims for future accepted formal rows.
